% !PRINTED 634
\weberpage{634}
%<<CONT
and consequently, so long as $r \leqq n$, by (4) with the notation (1):

\weq{12}{(n-r)a_r = a_r s_0 + a_{r-1}s_1 + \cdots + a_1 s_{r-1} + s_r,}

and if one takes the trace of $\theta^r f(\theta) = 0$:

\weq{13}{0 = a_n s_r + a_{n-1}s_{r+1} + \cdots + a_1 s_{r+n-1} + s_{r+n}.}

These are the \gesperrt{Newton}ian formulae (Bd.\ I, \S\,46).

If we denote by

\[
f'(t) = n t^{n-1} + (n-1)a_1 t^{n-2} + \cdots + 2a_{n-2}t + a_{n-1}
\]

the derived function of $f(t)$, multiply equation (12) by
$\theta^{n-r-1}$ and sum from $r=0$ to $r=n-1$, then it follows that:

\weq{14}{f'(\theta) = \eta_0 s_0 + \eta_1 s_1 + \cdots + \eta_{n-1}s_{n-1},}

and if, for any positive $k$, one multiplies equations (12) by
$\theta^{n-r+k-1}$ and takes in as many of equations (13) besides as are needed
until $n-r+k-1$ begins to become negative, there results:

\weq{15}{\theta^k f'(\theta) = \eta_0 s_k + \eta_1 s_{k+1} + \cdots + \eta_{n-1}s_{k+n-1}.}

To form the norm of $f'(\theta)$, one would have first to express, in these equations, the
basis $\eta_0, \eta_1, \ldots, \eta_{n-1}$ by means of the substitution (2) in terms of
$(1, \theta, \theta^2, \ldots, \theta^{n-1})$, then to form equation (15) for
$k = 0, 1, \ldots, n-1$ and to take the determinant of this system. Instead of
this one may take the determinant with respect to the $\eta$, and then multiply by the
substitution determinant:

\[
\begin{vmatrix}
a_{n-1}, & a_{n-2}, & \ldots, & a_1, & 1 \\
a_{n-2}, & a_{n-3}, & \ldots, & 1, & 0 \\
\cdots & \cdots & \cdots & \cdots & \cdots \\
1, & 0, & \ldots, & 0, & 0
\end{vmatrix}
= (-1)^{\frac{1}{2}n(n+1)}
% [sic] printed (-1)^{\frac{1}{2}n(n+1)}, but this determinant (1 on the anti-diagonal, 0 below it) equals (-1)^{\frac{1}{2}n(n-1)}; wrong sign for odd n
\]

In this way one obtains

\weq{16}{N[f'(\theta)] = (-1)^{\frac{1}{2}n(n+1)}
% [sic] printed (-1)^{\frac{1}{2}n(n+1)}, the factor is the substitution determinant, whose value is (-1)^{\frac{1}{2}n(n-1)}; wrong sign for odd n
\begin{vmatrix}
s_0, & s_1, & \ldots, & s_{n-1} \\
s_1, & s_2, & \ldots, & s_n \\
\cdots & \cdots & \cdots & \cdots \\
s_{n-1}, & s_n, & \ldots, & s_{2n-2}
\end{vmatrix}.}

But this determinant is, by the definition of the discriminant (\S\,173):

\[
\Delta(1, \theta, \theta^2, \ldots, \theta^{n-1}),
\]

and so we have:

\weq{17}{Nf'(\theta) = (-1)^{\frac{1}{2}n(n+1)}\Delta(1, \theta, \theta^2, \ldots, \theta^{n-1}).
% [sic] printed (-1)^{\frac{1}{2}n(n+1)}, the correct relation is Nf'(\theta) = (-1)^{\frac{1}{2}n(n-1)}\Delta(1, \theta, \ldots, \theta^{n-1}); wrong sign for odd n (e.g. n = 1)
}
